\documentclass[11pt]{article}
\usepackage[margin=1in]{geometry}
\usepackage{tabularx,booktabs,array,longtable}
\newcolumntype{Y}{>{\raggedright\arraybackslash}X}
\usepackage{amsmath,amssymb,bm}
\usepackage{xcolor}
\usepackage[hyphens]{url}
\usepackage{hyperref}
\sloppy
\emergencystretch=3em

\makeatletter
\newcommand{\code}[1]{{\def\_{\char`\_\discretionary{}{}{}}\texttt{#1}}}
\makeatother

\newcommand{\bU}{\mathbf{U}}
\newcommand{\bq}{\mathbf{q}}
\newcommand{\bv}{\mathbf{v}}
\newcommand{\bu}{\mathbf{u}}
\newcommand{\bB}{\mathbf{B}}
\newcommand{\bn}{\mathbf{n}}
\newcommand{\bF}{\mathbf{F}}
\newcommand{\Real}{\mathbb{R}}
\newcommand{\dd}{\,\mathrm{d}}
\newcommand{\Sbar}{\bar{S}}
\newcommand{\Cmax}{C_{\max}}
\newcommand{\Cmin}{C_{\min}}
\newcommand{\CR}{C_R}
\newcommand{\Cl}{C_l}
\newcommand{\Prt}{\mathrm{Pr}_t}
\newcommand{\dsgs}{\textsc{DynSGS}}

\title{The dynamic sub-grid-scale (\dsgs) model of JEXPRESSO\\
\large The equations exactly as implemented, and their correspondence with Dao \& Nazarov (2022)}
\author{}
\date{}

\begin{document}
\maketitle

\begin{abstract}
\noindent This note writes down, term by term, the residual-based artificial
viscosity (\dsgs) as it is coded in \code{src/kernel/physics/SGS.jl} and
applied in \code{src/kernel/operators/rhs.jl}: the discrete residual, its
normalization, the element and the nodal coefficient, the coefficient
assigned to every equation, and the viscous fluxes that are actually
assembled. Section~\ref{sec:table} tabulates, term by term, how each
ingredient is defined by Dao \& Nazarov (\emph{J.\ Sci.\ Comput.}\ 92:77,
2022, hereafter DN22) and how it is defined in JEXPRESSO. Nothing in this
note is a recommendation; it is a description of the code.
\end{abstract}

\tableofcontents

%======================================================================
\section{Setting and notation}
\label{sec:setting}

\paragraph{Unknowns.} The conserved state is a vector of $N_{eq}$ fields,
called \emph{slots}, in the following orders:
\begin{align}
\text{2D GLM-MHD (\code{DSGS\_MHD}, 9 slots):}\quad&
\bq = (\rho,\ \rho u,\ \rho v,\ E,\ \rho w,\ B_x,\ B_y,\ B_z,\ \psi), \label{eq:slots2d}\\
\text{1D MHD (\code{DSGS\_MHD}, 8 slots):}\quad&
\bq = (\rho,\ \rho u,\ \rho v,\ E,\ \rho w,\ B_x,\ B_y,\ B_z), \label{eq:slots1d}\\
\text{2D Euler, $\theta$ form (\code{DSGS}):}\quad&
\bq = (\rho,\ \rho u,\ \rho v,\ \rho\theta), \\
\text{2D Euler, energy form (\code{DSGS}):}\quad&
\bq = (\rho,\ \rho u,\ \rho v,\ \rho E), \\
\text{1D Euler (\code{DSGS}):}\quad&
\bq = (\rho,\ \rho u,\ \rho E).
\end{align}
$E$ is the total energy density. For MHD the pressure is
\begin{equation}
p = (\gamma-1)\Big(E - \tfrac12\rho|\bv|^2 - \tfrac12|\bB|^2 - \tfrac12\psi^2\Big),
\qquad \bv = (u,v,w),
\label{eq:pmhd}
\end{equation}
with $\psi\equiv 0$ in 1D; $\gamma$ is the deck's \code{:dsgs\_gamma}
(the case's $\gamma_{\rm mhd}$), not the air value of \code{PhysConst}.

\paragraph{Discretization.} Continuous Galerkin spectral elements of degree
$k$ on quadrilaterals (2D) or intervals (1D), $n_{gl}=k+1$ Legendre--Gauss--Lobatto
(LGL) points per direction, \emph{collocated} LGL quadrature (the flux is
interpolated at the nodes and integrated with the LGL rule; there is no
over-integration). The mass matrix is the LGL-lumped, diagonal matrix
$M = \mathrm{diag}(m_i)$; $M^{-1}$ is stored as the vector \code{Minv}. All
element contributions are assembled by direct stiffness summation (DSS).

For an element $K$, $\Delta_K$ is the mesh spacing stored in
\texttt{mesh.}$\Delta$\texttt{elem} (2D) or \texttt{mesh.}$\Delta$\texttt{x} (1D):
\begin{equation}
\Delta_K =
\begin{cases}
\min\big(\max_{c}x_c - \min_{c}x_c,\ \max_{c}y_c - \min_{c}y_c\big) & \text{2D, over the corner nodes $c$ of $K$},\\[2pt]
(x_{\max}-x_{\min})/N_{el} & \text{1D (uniform grid)}.
\end{cases}
\label{eq:DeltaK}
\end{equation}

\paragraph{Right-hand side.} At every evaluation of the semi-discrete
right-hand side (every Runge--Kutta stage), the inviscid weak form is built
element by element and assembled,
\begin{equation}
\mathrm{RHS}_i^{\,q} = -\int_\Omega \nabla\!\cdot\!\bF^{q}(\bq_h)\,\phi_i \dd x
\;+\; \int_\Omega S^{q}\,\phi_i \dd x ,
\qquad i = 1,\dots,N_h,
\label{eq:rhs}
\end{equation}
for each slot $q$ ($\bF^q$ its flux, $S^q$ its source), stored in
\code{params.RHS} \emph{before} the division by the mass matrix; the
element contributions $\mathrm{rhs}^{q}_{K,i}$ it is summed from (the same
integrals over one element $K$) are kept in \code{params.rhs\_el} and are
what the residual of Section~\ref{sec:residual} reads. The viscous
contribution of Section~\ref{sec:fluxes} is assembled separately into
\code{params.RHS\_visc}, added, and the sum is divided by $M$ to give
$\partial_t\bq_h$.

\paragraph{Per-slot multipliers.} The deck vector \texttt{:}$\mu$
$=(\mu_1,\dots,\mu_{N_{eq}})$ (called \code{visc\_coeff} in the code)
multiplies the coefficient of every slot. It is $1$ in all the MHD decks
except where noted, and it is written explicitly below.

%======================================================================
\section{The semi-discrete regularized system}
\label{sec:weak}

With $V_h$ the continuous nodal space, the code solves, for every slot $q$
and every test function $\phi_i\in V_h$,
\begin{equation}
\int_\Omega \partial_t q_h\,\phi_i \dd x
= \mathrm{RHS}_i^{\,q}
\;-\;\int_\Omega \bF_V^{q}(\bq_h)\cdot\nabla\phi_i \dd x ,
\label{eq:semidiscrete}
\end{equation}
with the mass matrix lumped. The viscous flux $\bF_V^{q}$ of each slot is
given in Section~\ref{sec:fluxes}. \textbf{No boundary integral}
$\int_{\partial\Omega}\bn\cdot\bF_V\,\phi_i$ is assembled: the diffusive
flux through the boundary is zero in the weak sense (homogeneous natural
condition), and Dirichlet data, where the case prescribes them, are imposed
strongly on the nodes after each stage.

%======================================================================
\section{The discrete residual}
\label{sec:residual}

\subsection{Time derivative: a step-cadenced history and a stage-consistent stencil}

\dsgs\ keeps its own three state snapshots, $\bq^{n}$, $\bq^{n-1}$ and
$\bq^{n-2}$ (\code{params.dsgs\_qn}, \code{dsgs\_qnm1}, \code{dsgs\_qnm2}),
advanced \emph{once per time step}: inside \code{rhs!}, when the stage time
satisfies $t - t_{\rm hist} \ge 0.999\,\Delta t$, i.e.\ at the first stage of
every step, the triple is rolled ($\bq^{n-2}\leftarrow\bq^{n-1}$,
$\bq^{n-1}\leftarrow\bq^{n}$, $\bq^{n}\leftarrow\bq_h$) and
$t_{\rm hist}\leftarrow t$. The Runge--Kutta stage snapshots
\code{qp.qnm1/qnm2} are \emph{not} used. $\Delta t$ is the deck's fixed time
step; there is no adaptive CFL time step.

At a stage with $\tau = t - t^n \in [0,\Delta t]$ and stage state $\bq_h$, the
time derivative is estimated as $w_1 q_i + w_2 q_{A,i} + w_3 q_{B,i}$ with
\begin{equation}
\begin{aligned}
\tau = 0:&\quad (q_A, q_B) = (q^{n-1}, q^{n-2}),\quad
w = \Big(\frac{3}{2h},\,-\frac{2}{h},\,\frac{1}{2h}\Big)\ \text{(BDF2)},\\
\tau > 0:&\quad (q_A, q_B) = (q^{n}, q^{n-1}),\quad
w = \Big(\frac{2\tau+h}{\tau(\tau+h)},\,-\frac{\tau+h}{\tau h},\,\frac{\tau}{h(\tau+h)}\Big),
\end{aligned}
\qquad h = \Delta t,
\label{eq:stencil}
\end{equation}
the second line being the three-point derivative at $\tau$ through
$(q(\tau), q^n, q^{n-1})$: second-order at every stage, equal to BDF2 at
$\tau = h$, weights summing to zero. (A fixed BDF2 on $(q(\tau), q^n, q^{n-1})$
reads $-\tfrac12\partial_t q$ at $\tau=0$ and drove the sensor to the cap
over every smooth moving structure.)

\subsection{The residual of slot $q$ at node $i$}

At every RHS evaluation (every stage), with $\bq_h$ the current stage
state,
\begin{equation}
\boxed{\;
R^{q,K}_i \;=\; \Big|\;w_1 q_i + w_2 q_{A,i} + w_3 q_{B,i}
\;-\; \frac{\mathrm{rhs}^{\,q}_{K,i}}{m^K_i}\;\Big| \;}
\qquad i \in K,
\label{eq:residual}
\end{equation}
where $\mathrm{rhs}^{q}_{K,i}$ is element $K$'s own weak inviscid right-hand
side at its node $i$ (fluxes and sources, before the direct stiffness
summation) and $m^K_i = \omega_i(\omega_j)J_{K,i}$ its lumped mass entry.
This is the \emph{element-wise} strong residual. With the lumped LGL mass
matrix the assembled rate $\mathrm{RHS}_i/m_i$ is exactly what the
integrator advances, so a residual built on it is the time-integration
error and nothing else; the element residual equals it at the interior
nodes of $K$ (summation-by-parts: $\mathrm{rhs}_{K,i}/m^K_i$ is the strong
nodal divergence) and differs from it at the interface nodes by the
mass-weighted jump of the flux divergence across the interface, $O(h^k)$
where the solution is smooth and $O(1/h)$ at a discontinuity. It is DN22's
$\frac1{m_i}\int|\mathrm{BDF}(q)+\nabla\!\cdot\! f(q)|\phi_i$ with the
absolute value inside the integral, evaluated with the LGL rule. The
element form takes $\max_{i\in K}R^{q,K}_i$; the nodal form the
mass-weighted average
\begin{equation}
R^{q}_i = \frac{\sum_{K\ni i} m^K_i\,R^{q,K}_i}{\sum_{K\ni i} m^K_i}
\qquad(\text{over the elements of the rank}).
\label{eq:Rnodal}
\end{equation}
When the case advances the total variables on top of a non-trivial
reference state $\bq_e$ (the hydrostatic atmosphere of the Euler $\theta$
cases, whose flux and source are the full ones) and the deck sets
\code{:dsgs\_reference => true}, $\mathrm{rhs}_{K,i}$
in~\eqref{eq:residual} is replaced by $\mathrm{rhs}_{K,i}(\bq_h) -
\mathrm{rhs}_{K,i}(\bq_e)$, the reference's element RHS being evaluated once:
at rest the element residual of the full equations is the interpolation
error of the hydrostatic balance at every interface, not an
under-resolution. Cases written on the perturbation (the well-balanced MHD
and shallow-water splits) have $\mathrm{rhs}_{K,i}(\bq_e) = 0$.
At the nodes of a Dirichlet boundary (a free-slip wall, a 1D end) the
boundary condition constrains the assembled rate while $\mathrm{rhs}_{K,i}$
carries the unconstrained tendency, so the element residual there is the
constraint force; it is set to zero at those nodes for every equation
(periodic and Laguerre edges are not constrained). A discontinuity reaching
a wall is still seen by the interior nodes of the same element.

\subsection{The residual set}

The maximum over equations (Section~\ref{sec:norm}) runs over the first
$N_{R}$ slots:
\begin{equation}
N_R = \min(N_{eq},8)\ \text{(MHD: $\rho,\rho u,\rho v,E,\rho w,B_x,B_y,B_z$; $\psi$ excluded)},\qquad
N_R = N_{eq}\ \text{(Euler)}.
\end{equation}

%======================================================================
\section{Normalization of the residual}
\label{sec:norm}

Each residual $R^q$ is divided by a scale $n^q$ with the units of $q$. Three
scopes exist, selected by \code{:dsgs\_norms} and by \code{:ldsgs\_nodal}.

\subsection{Domain (and rank) scope, element form}
\label{sec:norm-domain}

For each slot $q$, over the set $\mathcal{P}$ of \emph{element--node pairs}
$(K,i)$, $i\in K$ (so a node shared by $s$ elements is counted $s$ times),
\begin{equation}
\langle q\rangle = \frac{1}{|\mathcal{P}|}\sum_{(K,i)\in\mathcal{P}} q_i ,
\qquad
\Sbar^{q} = \max_{(K,i)\in\mathcal{P}} \big|q_i - \langle q\rangle\big| .
\label{eq:Sbar}
\end{equation}
With \code{:dsgs\_norms => "rank"} (the default) both are taken over the
rank's own elements without communication; with \code{"domain"} they are
reduced over all MPI ranks (the count $|\mathcal P|$, the sum, then the max).
On one rank the two coincide.

The spread $\Sbar^q$ is then bounded from below (``floored'': replaced by
$\max(\Sbar^q, a)$, so that it can never be smaller than $a$) at a
fraction $r = 10^{-3}$ of a physical scale of the slot built from the mean state
$\bar\rho = \max(|\langle\rho\rangle|,\epsilon)$,
$\bar p = (\gamma-1)\max\!\big(\langle E\rangle - \tfrac12(\langle\rho u\rangle^2+\langle\rho v\rangle^2+\langle\rho w\rangle^2)/\bar\rho - \tfrac12(\langle B_x\rangle^2+\langle B_y\rangle^2+\langle B_z\rangle^2),0\big)$,
$\bar c = \sqrt{\max(\gamma\bar p/\bar\rho,\epsilon)}$, $\epsilon=10^{-16}$:
\begin{equation}
n^{q} = \max\!\big(\Sbar^{q},\ r\,s^{q}\big) + \epsilon,
\qquad
s^{q} =
\begin{cases}
\bar\rho & q=\rho,\\
\bar\rho\,\bar c & q\in\{\rho u,\rho v,\rho w\},\\
\bar\rho\,\bar c^{\,2} & q = E,\\
\sqrt{\bar\rho}\,\bar c & q\in\{B_x,B_y,B_z\},
\end{cases}
\label{eq:floors}
\end{equation}
and $n^{\psi}=\Sbar^{\psi}+\epsilon$ (never used in the max).

\subsection{Element scope}
\label{sec:norm-element}

With \code{:dsgs\_norms => "element"} (\code{DSGS\_MHD} only), $n^q$ is
recomputed inside every element $K$ from the element's own $n_{gl}^d$
nodes: $\langle q\rangle_K$ and $\Sbar^q_K = \max_{i\in K}|q_i-\langle q\rangle_K|$,
and the lower bounds~\eqref{eq:floors} are applied with the \emph{element-mean}
state $(\bar\rho_K,\bar c_K)$ and the fraction $r_{\rm loc}$ =
\code{:dsgs\_local\_rel} (default $1$) in place of $10^{-3}$:
\begin{equation}
n^{q}_K = \max\!\big(\Sbar^{q}_K,\ r_{\rm loc}\,s^{q}_K\big) + \epsilon .
\label{eq:nelem}
\end{equation}
With $r_{\rm loc}=1$ a quiescent element measures its residual against the
local rate $\rho c$ per unit time.

\subsection{Nodal form (Dao--Nazarov normalization)}
\label{sec:norm-nodal}

With \code{:ldsgs\_nodal => true} the statistics are taken over the
\emph{unique} nodes $i=1,\dots,N_h$ (each counted once):
\begin{equation}
\langle q\rangle = \frac1{N_h}\sum_i q_i,\qquad
\Sbar^{q} = \max_i |q_i-\langle q\rangle|,\qquad
q_{\max}=\max_i q_i,\quad q_{\min}=\min_i q_i ,
\end{equation}
reduced over the ranks with \code{"domain"}, and $\Sbar^q$ is bounded from
below as in~\eqref{eq:floors} \emph{without} the trailing $+\epsilon$. For every node
$i$ the local range over its support $I(i)$ = the nodes of all elements that
contain $i$ (on the rank) is
\begin{equation}
\delta^{q}_i = \max_{j\in I(i)} q_j - \min_{j\in I(i)} q_j ,
\end{equation}
and the normalization is DN22 eq.~(4.7) with $\Cl$ = \code{:dsgs\_Cl}
(default $0$; DN22 use $0.4$):
\begin{equation}
\boxed{\;
n^{q}_i = \max\!\big(\Sbar^{q}, r s^{q}\big)\,
\Big(1 - \Cl\,\frac{\delta^{q}_i}{q_{\max}-q_{\min}}\Big)\;}
\qquad(\text{the bracket is } 1 \text{ if } q_{\max}-q_{\min}\le\epsilon).
\label{eq:n47}
\end{equation}

\subsection{The normalized residual}

\begin{align}
\text{element form:}\qquad&
\mathcal{R}_K = \max_{i\in K}\ \max_{q=1}^{N_R} \frac{R^{q}_i}{n^{q}}
\quad\big(n^{q}\to n^{q}_K \text{ in element scope}\big), \label{eq:RK}\\[4pt]
\text{nodal form:}\qquad&
\mathcal{R}_i = \max_{q=1}^{N_R}\ R^{q}_i\,\frac{n^{q}_i}{(n^{q}_i)^2+\epsilon},
\qquad \epsilon = 10^{-16}, \label{eq:Ri}
\end{align}
the latter being DN22 eq.~(4.8) with their division guard.

%======================================================================
\section{The kinematic viscosity}
\label{sec:nu}

\subsection{Wave speed}

At a node, with $\rho_i=\max(\rho,\epsilon)$, $\bv_i=\mathbf{m}_i/\rho_i$, $p_i$
from~\eqref{eq:pmhd} (clamped at $0$ in the nodal kernels and in 1D), and
$a^2=\gamma p_i/\rho_i$, $b^2=|\bB_i|^2/\rho_i$, $b_x^2 = B_{x,i}^2/\rho_i$:
\begin{align}
\text{2D:}\quad& \lambda_i = |\bv_i| + \sqrt{\max(a^2+b^2,0)}
\qquad\text{(an upper bound of the fast speed over all directions)},\label{eq:lam2d}\\
\text{1D:}\quad& \lambda_i = |\bv_i| + c_{f,i},\qquad
c_{f}^2 = \tfrac12\Big(a^2+b^2+\sqrt{(a^2+b^2)^2-4a^2 b_x^2}\Big)
\quad\text{(the exact fast speed along $x$)},\label{eq:lam1d}
\end{align}
where $|\bv_i| = \sqrt{u^2+v^2+w^2}$ with all three components.

\subsection{Element form (default, \code{:ldsgs\_nodal => false})}

For every element $K$, with
\begin{equation}
\Delta = \frac{\Delta_K}{n_{gl}} = \frac{\Delta_K}{k+1},\qquad
\lambda_K = \max_{i\in K}\lambda_i ,
\end{equation}
\begin{equation}
\boxed{\;
\begin{aligned}
\nu_{\rm res} &= \CR\,\Delta^2\,\mathcal{R}_K, &
\nu_{\max} &= \Cmax\,\Delta\,\lambda_K,\\
\nu_c &= \max\big(0,\ \min(\nu_{\max},\nu_{\rm res})\big), &
\nu_K &= \max\big(\nu_c,\ \Cmin\,\Delta\,\lambda_K\big)
\end{aligned}\;}
\label{eq:nuK}
\end{equation}
with $\CR$ = \code{:dsgs\_CR} (default $1$), $\Cmax$ = \code{:dsgs\_Cmax}
(default $0.5$), $\Cmin$ = \code{:dsgs\_Cmin} (default $0$). The last
operation is the background floor, a lower bound $\nu \ge \Cmin\Delta\lambda_K$
applied with a $\max$ (the counterpart of the cap $\nu \le \Cmax\Delta\lambda_K$
applied with a $\min$); with $\Cmin=0$ it is absent and
$\nu_K=\nu_c$. One value $\nu_K$ is used at every quadrature point of $K$;
$\nu$ is therefore piecewise constant and discontinuous across element
interfaces.

\subsection{Nodal form (\code{:ldsgs\_nodal => true})}

For every node $i$, with
\begin{equation}
h_i = \max_{K\ni i}\ \frac{\Delta_K}{k+1}
\qquad(\text{the element form's }\Delta\text{ of the elements containing } i),
\end{equation}
\begin{equation}
\boxed{\;
\nu_c = \max\big(0,\ \min(\Cmax\,h_i\,\lambda_i,\ \CR\,h_i^2\,\mathcal{R}_i)\big),\qquad
\nu_i = \max\big(\nu_c,\ \Cmin\,h_i\,\lambda_i\big)\;}
\label{eq:nui}
\end{equation}
(DN22 eqs.~(4.2), (4.10) plus the floor). $\nu_i$ is a continuous nodal
field in $V_h$: inside an element the coefficient at a quadrature point is
the nodal value at that (collocated) point, so the diffusive flux is
continuous across interfaces. The element means
$\bar\nu_K = n_{gl}^{-d}\sum_{i\in K}\nu_i$ are stored only for output.

%======================================================================
\section{Coefficients by equation}
\label{sec:byeq}

Let $\nu$ be $\nu_K$ (element form) or $\nu_i$ (nodal form), $\nu_c$ and
$\nu_{\rm fl}=\Cmin\Delta\lambda$ its residual and floor parts, and
$\mu_q$ the deck multiplier of slot $q$. Two forms exist, selected by
\code{:dsgs\_conserved}.

\subsection{Conserved form (\code{:dsgs\_conserved => true}; all MHD decks)}
\label{sec:conserved}

The case's \code{user\_primitives!} hands the assembly the conserved
variables themselves (or their departure from a reference state), and every
slot receives a kinematic coefficient:
\begin{center}
\renewcommand{\arraystretch}{1.4}
\begin{tabularx}{\textwidth}{l l Y}
\toprule
slot & coefficient & condition \\
\midrule
$\rho$ & $\mu_1\,\nu$ & \\
$\rho u,\ \rho v,\ \rho w$ & $\mu_{2,3,5}\,\nu$ & \\
$E$ & $\mu_4\,\nu$ & \code{:dsgs\_nazarov\_energy => false}, and always in 1D \\
$E$ & $\mu_4\,\max\!\big(\tfrac{\gamma(\gamma-1)}{\Prt}\,\nu_c,\ \nu_{\rm fl}\big)$ &
  \code{:dsgs\_nazarov\_energy => true}, 2D; acts on the thermal part only (Section~\ref{sec:split}) \\
$B_x,\ B_y,\ B_z$ & $\mu_{6,7,8}\,\nu$ & \\
$\psi$ & $\mu_9\,\nu$ & \\
\bottomrule
\end{tabularx}
\end{center}
\label{eq:cons}
The 1D kernels have no split energy flux. $\Prt$ = \code{:dsgs\_Prt} (default $0.7$; the MHD decks set $1$).

\subsection{Physical form (\code{:dsgs\_conserved => false})}
\label{sec:physical}

The case's \code{user\_primitives!} hands the assembly
$(\rho, u, v, T, w, B_x, B_y, B_z, \psi)$ with $T = p/\rho$, and
\begin{center}
\renewcommand{\arraystretch}{1.4}
\begin{tabularx}{\textwidth}{l l Y}
\toprule
slot & coefficient & condition \\
\midrule
$\rho$ & $\mu_1\,\nu$ & \\
$u,\ v,\ w$ & $\mu_{2,3,5}\,\varrho\,\nu$ & \\
$T$ & $\mu_4\,\varrho\,\nu\,\dfrac{\gamma}{(\gamma-1)\Prt}$ & \code{:dsgs\_nazarov\_energy => false}: Fourier law, $\kappa = c_p\varrho\nu/\Prt$ \\
$T$ & $\mu_4\,\varrho\,\nu/\Prt$ & \code{:dsgs\_nazarov\_energy => true}: DN22, $\kappa=\rho\nu/\mathrm{Pr}$ \\
$B_x,\ B_y,\ B_z,\ \psi$ & $\mu_{6,7,8,9}\,\nu$ & \\
\bottomrule
\end{tabularx}
\end{center}
\label{eq:phys}
where the density $\varrho$ that makes the coefficient dynamic is
\begin{equation}
\varrho =
\begin{cases}
\bar\rho_K = n_{gl}^{-d}\sum_{i\in K}\rho_i & \text{element form, \code{:dsgs\_nodal\_rho => false}},\\
\rho(\bm{x}_{q}) & \text{element form with \code{:dsgs\_nodal\_rho} (the quadrature point's density)},\\
\rho_i & \text{nodal form}.
\end{cases}
\end{equation}

\subsection{The split energy flux (\code{:dsgs\_conserved} with \code{:dsgs\_nazarov\_energy}, 2D)}
\label{sec:split}

The energy primitive is split by the case into a non-thermal and a thermal
part,
\begin{equation}
E = E_{\rm nth} + E_{\rm th},\qquad E_{\rm th} = \frac{p}{\gamma-1},\qquad
E_{\rm nth} = \tfrac12\rho|\bv|^2 + \tfrac12|\bB|^2 + \tfrac12\psi^2 ,
\end{equation}
slot 4 of \code{uprimitive} carrying $E_{\rm nth}$ (minus its reference
value) and the spare slot $N_{eq}+2$ carrying $E_{\rm th}$ (minus its
reference value). The energy flux assembled is then
\begin{equation}
\bF_V^{E} = \mu_2\,\nu\,\nabla E_{\rm nth} \;+\;
\mu_4\max\!\big(\tfrac{\gamma(\gamma-1)}{\Prt}\nu_c,\ \nu_{\rm fl}\big)\,\nabla E_{\rm th},
\label{eq:splitflux}
\end{equation}
so that the kinetic and magnetic energy removed by the $\nu\nabla(\rho\bv)$
and $\nu\nabla\bB$ Laplacians reappears in $E$ with the same coefficient,
while the thermal part conducts with
$\kappa = \rho\nu_c\,\gamma(\gamma-1)/\Prt$ on $\nabla(p/(\gamma-1))$,
i.e.\ $\kappa\nabla T$ with $\kappa=\rho\nu_c/\Prt$ to leading order
($T=p/\rho$, $\rho$ frozen), with the floor kept in full.

\subsection{The reference-weighted operator (\code{:dsgs\_ref\_weight}, 2D conserved form)}
\label{sec:refweight}

The case stores a positive weight $w(\bm{x})$ (the reference density
$\rho_e$ of \code{fluxEmergenceSon2025DSGS}) in the spare slot $N_{eq}+1$,
and hands slots $1$--$5$ as the relative departures $(q-q_e)/w$ and slots
$6$--$9$ as the absolute departures $q-q_e$. The assembled fluxes become
\begin{equation}
\bF_V^{q} = \mu_q\,\nu\,w\,\nabla\!\Big(\frac{q-q_e}{w}\Big),\ q=1..5,
\qquad
\bF_V^{q} = \mu_q\,\nu\,\nabla(q-q_e),\ q=6..9,
\end{equation}
with the energy split of~\eqref{eq:splitflux} applied to
$(E_{\rm nth}-E_{{\rm nth},e})/w$ and $(E_{\rm th}-E_{{\rm th},e})/w$.

%======================================================================
\section{The viscous fluxes assembled}
\label{sec:fluxes}

Write $\nu_q$ for the coefficient of slot $q$ from Section~\ref{sec:byeq}
(already including $\mu_q$ and, in the physical form, $\varrho$), $g_q$ for
the primitive of slot $q$ handed by the case, and $\eta = \nu_6$.

\subsection{2D (\code{\_expansion\_visc!}, ContGal, DSGS/DSGS\_MHD)}

At every LGL point of every element the flux vector $\bF_V^q$ of
\eqref{eq:semidiscrete} is
\begin{align}
q=1\ (\rho):\qquad& \bF_V = \nu_1\,\nabla g_1, \label{eq:fmass}\\[2pt]
q=2\ (x\text{-mom.}):\quad& \bF_V = \big(\tau_{xx},\ \tau_{xy}\big),\quad
\tau_{xx} = 2\nu_2\,\partial_x g_2 - \tfrac23\nu_2\,(\partial_x g_2+\partial_y g_3), \label{eq:fmomx}\\
&\phantom{\bF_V = \big(\tau_{xx},\ \tau_{xy}\big),\quad}
\tau_{xy} = \nu_2\,(\partial_y g_2+\partial_x g_3), \nonumber\\
q=3\ (y\text{-mom.}):\quad& \bF_V = \big(\tau_{xy},\ \tau_{yy}\big),\quad
\tau_{yy} = 2\nu_3\,\partial_y g_3 - \tfrac23\nu_3\,(\partial_x g_2+\partial_y g_3), \label{eq:fmomy}\\
&\phantom{\bF_V = \big(\tau_{xy},\ \tau_{yy}\big),\quad}
\tau_{xy} = \nu_3\,(\partial_y g_2+\partial_x g_3), \nonumber\\[2pt]
q=4\ (E):\qquad& \bF_V = \nu_4\,\nabla g_4
\;+\;\underbrace{\tau\cdot(g_2,g_3)}_{\text{physical form only}}
\label{eq:fE}\\
&\phantom{\bF_V = \nu_4\,\nabla g_4}\;+\;\underbrace{\eta\,\big(B_x\nabla g_6 + B_y\nabla g_7 + B_z\nabla g_8\big)}_{\text{physical form, \code{DSGS\_MHD} only}}, \nonumber\\[2pt]
q\ge5\ (\rho w,\ \bB,\ \psi):\qquad& \bF_V = \nu_q\,\nabla g_q. \label{eq:fscalar}
\end{align}
Remarks. (i) The stress uses the Stokes hypothesis
$\lambda=-\tfrac23\mu$; in the conserved form $g_2=\rho u$, $g_3=\rho v$
and the same expressions are applied to the momenta. (ii) The work term
$\tau\cdot\bu$ in~\eqref{eq:fE} uses the momentum coefficient $\nu_2$
recomputed at the point and the primitives $g_2,g_3$ as velocities; it is
skipped in the $\theta$ form and in the conserved form. (iii) The resistive
work term uses the component-Laplacian form $\eta\,\bB\cdot\nabla\bB$, the
energy counterpart of the induction term~\eqref{eq:fscalar}; it is not the
$\eta\bB\cdot(\nabla\bB-\nabla\bB^{\!\top})$ of DN22 eq.~(4.4). (iv) With
the split energy flux, $\nu_4\nabla g_4$ in~\eqref{eq:fE} is replaced
by~\eqref{eq:splitflux}. (v) The induction slots receive a scalar Laplacian
per component, $\eta\nabla B_x$, $\eta\nabla B_y$, $\eta\nabla B_z$, not the
antisymmetric tensor $\eta(\nabla\bB-\nabla\bB^{\!\top})$.

The weak form is assembled as
\begin{equation}
\mathrm{RHS}^{\rm visc}_{i}
\mathrel{-}= \sum_{\text{LGL pts}} \omega_k\omega_l J_{kl}\;
\nabla\phi_i(\bm{x}_{kl})\cdot\bF_V(\bm{x}_{kl}),
\end{equation}
with $\nabla\phi_i$ built from the reference derivatives and the metric
terms at the point, exactly as for any other viscous model of the code.

\subsection{1D (\code{\_expansion\_visc!}, DSGS, NSD\_1D)}

Every slot receives a scalar diffusion,
\begin{equation}
\bF_V^{q} = \nu_q\,\partial_x g_q ,\qquad
\mathrm{RHS}^{\rm visc}_{i} \mathrel{-}= \sum_{k}\omega_k J_k\,(\partial_x\phi_i)_k\,\nu_{q,k}\,(\partial_x g_q)_k ,
\end{equation}
with $\nu_{q,k}=\nu_{q,K}$ (element form) or the nodal value $\nu_{q,i(k)}$
(nodal form). There is no deviatoric stress, no $\tau\cdot u$ and no
resistive work term in 1D; the conserved form on
$(\rho,\rho u,\rho v,E,\rho w,B_x,B_y,B_z)$ is exactly conservative.

%======================================================================
\section{The Euler kernels (\code{DSGS}) and the shallow-water kernel (\code{DSGS\_SW})}
\label{sec:euler}

These predate the MHD kernel and keep their own hard-coded constants.

\subsection{2D, $\theta$ form (\code{:energy\_equation => "theta"})}

$\bq=(\rho,\rho u,\rho v,\rho\theta)$, $\CR=1$, $\Cmax=0.5$ (not deck
inputs), $p = C_0(\rho\theta)^\gamma$ with $C_0,\gamma$ from
\code{PhysConst}, $c=\sqrt{\gamma p/\rho}$. Domain scope of
Section~\ref{sec:norm-domain} with the lower bounds replaced by $n^q = \Sbar^q+\epsilon$
for $\rho$ and $\rho\theta$ and $n^q=\max(\Sbar^q, 10^{-3}\bar\rho\bar c)+\epsilon$
for the momenta. Then per element,
$\nu_K$ as in~\eqref{eq:nuK} with $\Cmin=0$, $\lambda_i=|\bv_i|+c_i$, and the coefficients
\begin{equation}
\rho:\ 0,\qquad u,v:\ \mu_{2,3}\,\bar\rho_K\,\nu_K,\qquad
\theta:\ \mu_4\,\frac{\mathrm{Pr}}{\gamma-1}\,\bar\rho_K\,\nu_K ,\qquad
q\ge5:\ \mu_q\,\nu_K ,
\end{equation}
with $\mathrm{Pr}$ = \code{:Pr}, the artificial Prandtl number of
Nazarov \& Hoffman (default $0.1$, the value of every DSGS deck). The primitives are
$(\rho,u,v,\theta)$; slots $q\ge5$ are passive tracers
(\code{CompEuler/thetaTracers}), transported un-weighted and diffused with the
kinematic $\nu_K$; they enter neither the residual nor the normalization.

\subsection{2D, energy form (\code{:energy\_equation => "energy"}; Nazarov \& Hoffman 2013)}

$\bq=(\rho,\rho u,\rho v,\rho E)$, $\CR=1$, $\Cmax=0.5$. The momentum
residual and spread are taken as vectors:
$R^{m}_i = \sqrt{(R^{\rho u}_i)^2+(R^{\rho v}_i)^2}$,
$\Sbar^{m}=\max_i\sqrt{(\rho u_i-\langle\rho u\rangle)^2+(\rho v_i-\langle\rho v\rangle)^2}$,
bounded from below as in~\eqref{eq:floors} for $\rho$, $m$, $E$. With
$T_i = \max(E_i/\rho_i - \tfrac12|\bv_i|^2,0)$ and $\lambda_i=|\bv_i|+\sqrt{\gamma T_i}$,
per element ($h=\Delta_K/n_{gl}$)
\begin{equation}
\mu_{\rm res} = \CR\,h^2\,\Sbar^{\rho}\,\mathcal{R}_K,\qquad
\mu_{\rm cap} = \Cmax\,h\,\max_{i\in K}\rho_i\,\max_{i\in K}\lambda_i,\qquad
\mu_K = \max(0,\min(\mu_{\rm cap},\mu_{\rm res})),
\end{equation}
which is a \emph{dynamic} viscosity, and the coefficients
\begin{equation}
\rho:\ \mu_1\,\frac{\mu_K}{\max_{i\in K}\rho_i},\qquad
u,v:\ \mu_{2,3}\,\mu_K,\qquad
T:\ \mu_4\,\frac{\mathrm{Pr}}{\gamma-1}\,\mu_K ,\qquad
q\ge5:\ \mu_q\,\frac{\mu_K}{\max_{i\in K}\rho_i} ,
\end{equation}
on the primitives $(\rho,u,v,T)$ with $T = E/\rho-\tfrac12|\bv|^2$; the
$\tau\cdot\bu$ term of~\eqref{eq:fE} is included.

\subsection{3D, $\theta$ form (\code{CompEuler/3d}, the LES cases)}
\label{sec:euler3d}

$\bq=(\rho,\rho u,\rho v,\rho w,\rho\theta)$: the $\theta$ kernel above with the
node loop over the three directions and the lumped mass entry
$m^K_i=\omega_i\omega_j\omega_k J_{K,ijk}$. All five residuals enter the max;
the momentum floor $10^{-3}\bar\rho\bar c$ is applied to all three momenta
(the atmospheric cases start globally at rest, so those spreads start at
exactly zero); $\Delta=\Delta_K/n_{gl}$, $\lambda_i=|\bv_i|+c_i$ with
$c=\sqrt{\gamma p/\rho}$ from $p=C_0(\rho\theta)^\gamma$. The primitives are
$(\rho,u,v,w,\theta)$ and the slot coefficients
\begin{equation}
\rho:\ 0,\qquad u,v,w:\ \mu_{2,3,4}\,\bar\rho_K\,\nu_K,\qquad
\theta:\ \mu_5\,\frac{\mathrm{Pr}}{\gamma-1}\,\bar\rho_K\,\nu_K,\qquad
q\ge6:\ \mu_q\,\nu_K .
\end{equation}
The residual is zeroed on the nodes of the non-periodic boundary
\emph{faces} (\code{mesh.poin\_in\_bdy\_face}), the 3D counterpart of the
boundary-edge rule of Section~\ref{sec:residual}. The viscous operator is the
per-slot Laplacian $\nabla\cdot(\mu_q\nabla q)$, not the 2D path's stress
form. Only the $\theta$ form exists in 3D: \code{:energy\_equation => "energy"}
with \code{DSGS()} in 3D raises rather than building a $\theta$-form
coefficient from a total-energy state, and there is no 3D nodal form.

Before September 2026 there was no 3D kernel at all, and the omission was
silent: \code{params.sgs} is \code{nothing} for every model but Smagorinsky
and Vreman, so a 3D case with \code{:visc\_model => DSGS()} reached the
constant-coefficient branch of \code{\_viscous\_rhs\_el\_3d!} and ran the deck's
per-equation multipliers $\mu_q$ as plain Laplacian coefficients in m$^2$/s.

\subsection{2D nodal form of the Euler kernel}

With \code{:ldsgs\_nodal => true}, $\nu_i$ from~\eqref{eq:nui} with the
per-slot normalization~\eqref{eq:n47} (every momentum component on its own),
$\CR,\Cmax,\Cmin,\Cl$ from the deck, $\lambda_i=|\bv_i|+c_i$ ($c$ from the
$\theta$ or the energy equation of state), and the slot coefficients
$\theta$ form: $(0,\ \rho_i\nu_i,\ \rho_i\nu_i,\ \tfrac{\mathrm{Pr}}{\gamma-1}\rho_i\nu_i)$;
energy form: $(\nu_i,\ \rho_i\nu_i,\ \rho_i\nu_i,\ \tfrac{\mathrm{Pr}}{\gamma-1}\rho_i\nu_i)$,
passive tracers $\nu_i$, each times $\mu_q$.

\subsection{1D Euler}

$\bq=(\rho,\rho u,\rho E)$, $\gamma=1.4$, $\CR=1$, $\Cmax=0.5$,
$n^q=\Sbar^q+\epsilon$ (no physical lower bounds), $\Delta=\Delta_K/n_{gl}$,
$\lambda_i=|u_i|+\sqrt{\gamma(\gamma-1)e_i}$ with
$e_i=\max(E_i/\rho_i-\tfrac12u_i^2,0)$, and the \emph{same} $\nu_K$ on all
three slots, times $\mu_q$.

\subsection{2D shallow water (\code{DSGS\_SW})}

$\bq=(H,\,Hu,\,Hv)$, $\bq_e = (H_e,0,0)$ the lake at rest of the case,
$\CR,\Cmax,\Cmin,\Cl$ from the deck. The statistics of
Section~\ref{sec:norm} are taken on the departure $\delta\bq = \bq-\bq_e$
(the depth itself is cone-shaped over an island), with the lower bounds
$n^{H} = \max(\Sbar^{H}, 10^{-3}\bar H)$,
$n^{Hu} = n^{Hv} = \max(\Sbar, 10^{-3}\bar H\sqrt{g\bar H})$, $\bar H$ the
mean depth; the residual set is the three equations. The wave speed is
$\lambda_i = |\bv_i| + \sqrt{g\max(H_i,0)}$ with
$\bv_i = (Hu, Hv)_i/\max(H_i, h_{min})$, $h_{min}$ = \code{:dsgs\_swe\_hmin}
(the case's wet/dry threshold) and $g$ = \code{:dsgs\_swe\_g}. Then $\nu_K$
as in~\eqref{eq:nuK} (element form) or $\nu_i$ as in~\eqref{eq:nui} (nodal
form), and the \emph{same} kinematic $\nu$ on the three slots, times
$\mu_q$, applied on the primitives $(H-H_e,\,Hu,\,Hv)$ through the 2D
fluxes~\eqref{eq:fmass}--\eqref{eq:fmomy} (no energy slot).

%======================================================================
\section{Output of the coefficients}

\texttt{params.}$\mu$\texttt{\_dsgs[K,q]} holds $\nu_q$ per element (element form) or
the element mean of the nodal $\nu_{q,i}$ (nodal form);
\texttt{params.}$\mu$\texttt{\_dsgs\_pnode[i,q]} holds the value written to VTK/PNG:
the nodal field itself in the nodal form, and in the element form the
element value broadcast to the nodes (a shared node receives the value of
the last element that owns it). The writer emits one field
\code{mu\_dsgs\_<slots>} per distinct coefficient.

%======================================================================
\section{Deck keys}

\begin{center}
\small
\begin{tabularx}{\textwidth}{l l Y}
\toprule
key & default & meaning \\
\midrule
\code{:visc\_model} & -- & \code{DSGS\_MHD()} or \code{DSGS()} \\
\texttt{:}$\mu$ & -- & per-slot multipliers $\mu_q$ \\
\code{:dsgs\_CR} & 1.0 & $\CR$ (residual viscosity) \\
\code{:dsgs\_Cmax} & 0.5 & $\Cmax$ (first-order viscosity) \\
\code{:dsgs\_Cmin} & 0.0 & $\Cmin$ (background floor; not in DN22) \\
\code{:dsgs\_Cl} & 0.0 & $\Cl$ of eq.~\eqref{eq:n47} (nodal form) \\
\code{:dsgs\_gamma} & 5/3 & $\gamma$ of the MHD kernel \\
\code{:dsgs\_Prt} & 0.7 & $\Prt$ of the MHD energy slot \\
\code{:dsgs\_norms} & \code{"rank"} & scope of $\langle q\rangle,\Sbar^q$: \code{"rank"} (no reductions), \code{"domain"}, \code{"element"} \\
\code{:dsgs\_local\_rel} & 1.0 & $r_{\rm loc}$ of eq.~\eqref{eq:nelem} \\
\code{:ldsgs\_nodal} & false & nodal form (\code{true}) or element form \\
\code{:dsgs\_conserved} & false & conserved form (Section~\ref{sec:conserved}) \\
\code{:dsgs\_nazarov\_energy} & false & $\kappa=\rho\nu/\Prt$ instead of $c_p\rho\nu/\Prt$; split energy flux in the conserved form \\
\code{:dsgs\_nodal\_rho} & false & $\varrho=\rho(\bm{x}_q)$ in the physical form \\
\code{:dsgs\_ref\_weight} & false & reference-weighted operator (Section~\ref{sec:refweight}) \\
\code{:Pr} & 0.1 & artificial $\mathrm{Pr}$ of the Euler kernels ($\kappa = \tfrac{\mathrm{Pr}}{\gamma-1}\mu$) \\
\bottomrule
\end{tabularx}
\end{center}

%======================================================================
\section{Dao \& Nazarov (2022) versus JEXPRESSO, term by term}
\label{sec:table}

The table compares DN22 with the \code{DSGS\_MHD} kernel. ``Element'' and
``nodal'' refer to \code{:ldsgs\_nodal => false / true}. Where the two are
identical the entry says so.

\begin{small}
\begin{longtable}{@{}p{0.18\textwidth} p{0.37\textwidth} p{0.39\textwidth}@{}}
\toprule
term & Dao \& Nazarov 2022 & JEXPRESSO \\
\midrule
\endfirsthead
\toprule
term & Dao \& Nazarov 2022 & JEXPRESSO \\
\midrule
\endhead
\bottomrule
\endfoot

Discretization &
Continuous Lagrange $P_k$ ($P_1$, $P_3$) on simplices, FEniCS, inner products integrated exactly &
Continuous LGL spectral elements of degree $k$ on quadrilaterals / intervals, collocated LGL quadrature (inexact) \\

Mass matrix &
Consistent $M$ in the solve $MK_l = F$; lumped $m_i=\int\phi_i$ in the residual only &
Lumped (diagonal) $M$ everywhere; the residual uses the element mass entries $m^K_i$ \\

Unknowns in the residual &
$\rho,\ m_1..m_d,\ E,\ B_1..B_d$ (eq.~4.8); $\Psi_l$ of the GLM system not included &
$\rho,\ \rho u,\ \rho v,\ E,\ \rho w,\ B_x,\ B_y,\ B_z$; $\psi$ excluded. Identical set \\

Time derivative in the residual &
$\mathrm{BDF2}(q)^n = (3q^n-4q^{n-1}+q^{n-2})/(2h)$, $h$ the time step &
BDF2 at the first stage of a step; at a later stage the second-order three-point derivative through the stage state, $q^n$ and $q^{n-1}$, eq.~\eqref{eq:stencil}; the snapshots are once-per-step \\

Residual $R(q)$ &
Lumped $L^2$ projection of the absolute value: $R_i = \frac1{m_i}\int|\mathrm{BDF2}+\nabla\!\cdot f|\,\phi_i$ (eq.~4.9) &
Element-wise: $R^K_i = |\partial_t q_i - \mathrm{rhs}_{K,i}/m^K_i|$ with the element's own weak RHS and lumped mass entry, eq.~\eqref{eq:residual}; nodal form: mass-weighted average over the elements of the node, eq.~\eqref{eq:Rnodal}; the absolute value is inside the average as in DN22 \\

When $\nu$ is computed &
Once per time step from $U^n$; frozen over the RK stages &
At every RHS evaluation (every RK stage) from the stage state, with the step-cadenced history \\

Time step &
Adaptive, $\tau_n = \mathrm{CFL}\,\min h_h/\max\lambda_{\max}$, CFL $=0.3$, classical RK4 &
Fixed $\Delta t$ from the deck; SSPRK53 or CarpenterKennedy2N54 (deck choice) \\

Mesh function $h$ &
$h_h\in Q_h$, the $L^2$ projection of $h_K/k$, $h_K$ the circumradius of $K$ (eq.~2.2) &
Element: $\Delta=\Delta_K/(k+1)$ with $\Delta_K$ the min corner extent~\eqref{eq:DeltaK}. Nodal: $h_i=\max_{K\ni i}\Delta_K/(k+1)$, the same length. (DN22's $h_K/k$ with the circumradius is $\Delta_K/(\sqrt2 k)$ on a square, within 12\,\% of $\Delta_K/(k+1)$ at $k=4$; $\Delta_K/k$ gave a start-up coefficient 1.56$\times$ the element form's on the rising-bubble case and exceeded the explicit viscous limit) \\

Wave speed $\lambda_{\max}$ &
$\max_{i=1..8}|\lambda_i|$ of the eigenvalues~(3.6) along a direction $e$: $|u\!\cdot\!e| + c_f(e)$ &
2D: $|\bv|+\sqrt{a^2+|\bB|^2/\rho}$, the bound of $c_f$ over all directions, with the full 3-component $|\bv|$. 1D: $|\bv|+c_f(x)$, exact. Element form takes the max over the element's nodes \\

First-order viscosity &
$\nu^L_i = \Cmax h_i\lambda_{\max,i}$, $\Cmax=0.5$ (\S4.2) &
Element: $\Cmax\Delta\lambda_K$. Nodal: identical, $\Cmax h_i\lambda_i$. \code{:dsgs\_Cmax}, default $0.5$ \\

Global spread $\Sbar(w)$ &
$\|w-\mathrm{mean}(w)\|_{L^\infty(\Omega)}$ &
Same, with the mean and max over the nodes (nodal form: unique nodes; element form: element--node pairs, shared nodes counted once per element); over the rank only with \code{:dsgs\_norms => "rank"} \\

Lower bounds on $\Sbar$ &
None; the ratio uses $R\,n/(n^2+\epsilon)$, $\epsilon\approx2.2\times10^{-16}$ &
$\Sbar^q \leftarrow \max(\Sbar^q, 10^{-3}s^q)$ with $s^q$ the slot's physical scale~\eqref{eq:floors} from the mean state; then the same guard in the nodal form, $R/(n+\epsilon)$ in the element form \\

Local-jump normalization &
$n(w)_i = \Sbar(w)\big(1-\Cl\,\frac{\max_{I(i)}w-\min_{I(i)}w}{\max w-\min w}\big)$, eq.~4.7, $\Cl=0.4$, $I(i)$ the neighbours of node $i$ &
Nodal: identical, eq.~\eqref{eq:n47}, $I(i)$ = the nodes of the elements containing $i$, \code{:dsgs\_Cl} (default $0$). Element: not applied ($\Cl=0$); \code{"element"} scope uses the element spread~\eqref{eq:nelem} instead \\

Normalized residual &
$R^n_h = \max_q R(q)/n(q)$ at every node, eq.~4.8 &
Nodal: identical, eq.~\eqref{eq:Ri}. Element: $\max$ over the element's nodes and the equations, eq.~\eqref{eq:RK} \\

High-order viscosity &
$\nu^H_i = \CR h_i^2 R^n_i$, $\CR=1$ (eq.~4.10) &
Element: $\CR\Delta^2\mathcal R_K$. Nodal: identical. \code{:dsgs\_CR}, default $1$ \\

Combination &
$\nu_i = \min(\nu^L_i,\nu^H_i)$ (eq.~4.10) &
$\max(0,\min(\nu_{\max},\nu_{\rm res}))$, then $\max(\cdot,\Cmin\Delta\lambda)$; with $\Cmin=0$ identical \\

Background floor &
None &
$\Cmin\,h\,\lambda$ (\code{:dsgs\_Cmin}), $0$ by default; the Brio--Wu deck uses $0.03$ \\

Where $\nu$ lives &
Nodal field in $Q_h$ (continuous) &
Nodal: identical. Element: one value per element (discontinuous), the default \\

Mass equation &
$\nu_h\nabla\rho_h$ (eq.~4.4) &
$\mu_1\nu\nabla\rho$; identical with $\mu_1=1$ (the MHD decks) \\

Dynamic viscosity &
$\mu_h = \rho_h\nu_h$ (\S4.4), nodal $\rho$ &
Physical form: $\varrho\nu$ with $\varrho$ = element-mean $\bar\rho_K$ (default), nodal $\rho$ (\code{:dsgs\_nodal\_rho} or nodal form). Conserved form: $\nu\nabla(\rho\bv)$, no explicit $\rho$ \\

Stress tensor &
$\tau = \mu(\nabla u+\nabla u^{\!\top}) + \lambda\nabla\!\cdot\!u\,I$ with $\lambda=0$ &
$\tau = \mu(\nabla u+\nabla u^{\!\top}) - \tfrac23\mu\nabla\!\cdot\!u\,I$ (Stokes hypothesis), eqs.~\eqref{eq:fmomx}--\eqref{eq:fmomy}; applied to $(u,v)$ in the physical form and to $(\rho u,\rho v)$ in the conserved form. The third component $\rho w$ receives a scalar Laplacian. 1D: scalar Laplacian on every slot \\

Heat conduction &
$\kappa_h = \mu_h/\mathrm{Pr}$ on $\nabla T_h$, $T=p/\rho$ &
Physical form: $\varrho\nu\,\gamma/((\gamma-1)\Prt)$ on $\nabla(p/\rho)$ by default (Fourier, $c_p$), or $\varrho\nu/\Prt$ with \code{:dsgs\_nazarov\_energy} (identical to DN22). Conserved form: $\nu\nabla E$ by default; with \code{:dsgs\_nazarov\_energy} the split flux~\eqref{eq:splitflux}, $\kappa=\rho\nu_c/\Prt$ to leading order on the thermal part with the floor kept \\

Viscous work in $E$ &
$u_h\cdot\tau_h$ &
Physical form: $\tau\cdot(g_2,g_3)$, same. Conserved form: not added (the $\nu\nabla E$ Laplacian carries it) \\

Resistive work in $E$ &
$\eta_h B_h\cdot(\nabla B_h-\nabla B_h^{\!\top})$ &
Physical form: $\eta\,\bB\cdot\nabla\bB$ (component-Laplacian form, matching the induction term below). Conserved form: not added \\

Induction equation &
$\eta_h(\nabla B_h-\nabla B_h^{\!\top})$, $\eta_h=\nu_h$ &
$\eta\nabla B_x$, $\eta\nabla B_y$, $\eta\nabla B_z$ (a scalar Laplacian per component), $\eta=\mu_{6,7,8}\nu$ \\

Divergence cleaning &
Projection, pseudo-time, or hyperbolic GLM ($\Psi_l$, $c_h=\max\lambda$, damping $c_r=0.3$ in 2D) &
2D: Dedner mixed GLM with $\psi$ as slot 9 (case files); $\psi$ is diffused with $\mu_9\nu$ and excluded from the residual. 1D: $B_x$ constant, no cleaning \\

Boundary term of the viscous form &
$(n\cdot F_V, V_h)_\Gamma$ kept (Neumann as a boundary integral); Dirichlet and slip strongly &
Boundary integral not assembled (zero diffusive flux); Dirichlet strongly on the nodes \\

Bulk viscosity &
$\lambda_h = 0$ &
$-\tfrac23\mu$ (through the deviatoric stress) \\

Parameters used &
$\Cmax=0.5$, $\CR=1$, $\Cl=0.4$, CFL $0.3$, $\mathrm{Pr}$ (unstated in \S5) &
Brio--Wu: $\Cmax=0.5$, $\CR=1$, $\Cl=0.4$ (nodal), $\Cmin=0.03$, $\Prt=1$, conserved form, element form by default \\

\end{longtable}
\end{small}

\end{document}
